% !TeX program = xelatex
% UTF-8 transcription of all ten source pages; source wording is retained.
% Original images are images/pic01.jpg through images/pic10.jpg.
% Standalone document, or input into an existing XePersian document.
\ifdefined\settextfont
  \def\MyocarditisStandalone{0}
\else
  \def\MyocarditisStandalone{1}
  \documentclass[12pt,a4paper]{article}
  \usepackage[margin=2cm]{geometry}
  \usepackage{amsmath,amssymb,enumitem}
  \usepackage{xepersian}
  \IfFontExistsTF{Amiri}{\settextfont{Amiri}}{\settextfont{FreeSerif}}
  \setlatintextfont{FreeSerif}
\fi
\providecommand{\en}[1]{\lr{#1}}
\newcommand{\myopage}[2]{\clearpage
  \section*{موضوع: براونوالد ـ ۵۵ ـ میوکاردیت ـ صفحهٔ #1}
  {\footnotesize تاریخ: \en{/ /}\quad\en{#2}\par}}
\ifnum\MyocarditisStandalone=1 \begin{document}\fi

\myopage{۱}{pic01.jpg — 20261012\_185801976.jpg}
\begin{itemize}
\item میوکاردیت: هر التهابی در میوکارد قلب.
\item ایسکمی، تروما، کاردیومیوپاتی.
\item تعریف جزئی‌تر: التهاب میوکارد قلب به دنبال مواجهه با آنتی‌ژن خارجی (ویروسی، باکتری، قارچ، انگل) یا داخلی (خودایمنی).
\item شایع‌ترین علت شناخته‌شدهٔ میوکاردیت: عفونت ویروسی.
\item شایع‌ترین ویروس‌ها:
  \begin{itemize}
  \item پاروویروس \en{$B_{19}$} (\en{$B_{19}V$}).
  \item هرپس ویروس انسانی ۶ (\en{HHV6}).
  \end{itemize}
\item اگر پاسخ میزبان به ویروس بیش از حد باشد، منجر به مرگ میوسیت‌ها می‌شود که به \en{DCM}، ریمودلینگ قلبی، \en{HF} و مرگ می‌انجامد.
\item در برخی موارد ژنوم ویروس در بیوپسی اندومیوکاردیال (\en{EMB}) در بیماران \en{DCM} یافت می‌شود.
\end{itemize}
\subsection*{سندرم بالینی در میوکاردیت}
% Heading continues on page 2.

\myopage{۲}{pic02.jpg — 20261012\_185808973.jpg}
\begin{itemize}
\item در یک سر طیف ممکن است بیمار بدون علامت باشد و ما در \en{EKG} یا \en{Echo} به بیماری پی ببریم و در سر دیگر بیمار شدیداً علامت‌دار است.
\end{itemize}
\subsection*{یافته‌ها در میوکاردیت}
\begin{latin}
Chest Pain, Cardiac Dysfunction, Arrhythmias, HF, hemodynamic collapse
\end{latin}
\begin{itemize}
\item در میوکاردیت‌های ویروسی تغییرات \en{EKG} و \en{Echo} موقت ممکن است دیده شود.
\item \en{Chest Pain} در میوکاردیت مانند آنژین است و می‌تواند همراه با تغییرات \en{EKG} باشد: مثل \en{STE}.
\item میوکاردیت به دو گروه:
  \begin{itemize}
  \item زودرس: شروع حاد و ناگهانی.
  \item اواخر سیر: شروع مزمن همراه با \en{DCM} و \en{HF}.
  \end{itemize}
\item میوکاردیت عامل \en{10--50\%} موارد \en{DCM} ایدیوپاتیک است.
\item \en{HF} در میوکاردیت هم سیستولیک و هم دیاستولیک می‌تواند باشد.
\end{itemize}

\myopage{۳}{pic03.jpg — 20261012\_185814121.jpg}
\begin{itemize}
\item تظاهرات بالینی میوکاردیت وابسته به اتیولوژی آن است.
  \begin{itemize}
  \item \en{$B_{19}V$}: همراه با \en{CP}.
  \item جاینت سل میوکاردیت (\en{GCM}): آریتمی بطنی، بلاک.
  \item سارکوئیدوز: تنگی نفس و آدنوپاتی مدیاستن در \en{CXR}؛ بلاک.
  \item حساسیت به دارو: خارش، راش ماکولوپاپولار، \en{$\uparrow$} ائوزینوفیل.
  \end{itemize}
\item در بالین افتراق بین یک \en{DCM} غیرالتهابی و یک \en{DCM} التهابی مزمن نیازمند انجام \en{EMB} است.
\end{itemize}
\subsection*{رویکرد تشخیصی}
بر اساس کرایتریا دالاس، تشخیص کلاسیک نیازمند بیوپسی است.
\subsection*{تظاهرات بالینی در میوکاردیت حاد}
\begin{itemize}
\item خستگی (۸۲\%)، تنگی نفس فعالیتی (۸۱\%)، آریتمی (۵۵\%)، تپش قلب (۴۹\%)، \en{CP} در حالت استراحت (۲۶\%).
\item افتراق از ایسکمی دشوار است چون \en{Trop} بالا می‌رود و تغییرات \en{EKG} هم داریم. تغییرات \en{Echo} هم می‌تواند مشابه باشد.
\item پرودروم بیماری ویروسی (سرماخوردگی) در \en{20--80\%} موارد وجود دارد.
\end{itemize}

\myopage{۴}{pic04.jpg — 20261012\_185821500.jpg}
\subsection*{تظاهرات بالینی میوکاردیت فولمیننت}
\begin{itemize}
\item ظرف ۲ هفته از شروع بیماری ویروسی.
\item اختلال همودینامیک و افت \en{BP}: اغلب نیازمند حمایت مکانیکی.
\item در \en{Echo}: هایپوفانکشن منتشر؛ ضخیم شدن دیوارهٔ بطنی (به دلیل ادم میوکارد).
\item در \en{EMB} التهاب منتشر در هر نقطه‌ای دیده می‌شود.
\item ۹۳\% زنده می‌مانند و نیاز به پیوند قلب پیدا نمی‌کنند.
\item باید به صورت تهاجمی آن‌ها را درمان کرد.
\end{itemize}
\subsection*{تظاهرات بالینی \en{GCM}}
\begin{itemize}
\item در ابتدا از سایر میوکاردیت‌ها افتراق داده نمی‌شود.
\item تظاهر با بلاک، آریتمی، \en{HF}.
\item علی‌رغم درمان این موارد بهبود پیدا نمی‌کنند.
\item در صورت عدم درمان بیمار ظرف ۶ ماه می‌میرد.
\item درمان: \en{high dose multiagent immunosuppression}.
\item قطع زودهنگام درمان منجر به عود می‌شود.
\item بیوپسی:
\end{itemize}
\begin{latin}
Distinctive pattern of giant cells, Active inflammation and Scar tissue.
\end{latin}
\begin{itemize}
\item تنها درمان مؤثر نهایی پیوند قلب است.
\item همراه با سایر اختلالات خودایمنی مثل کرون و تایموما همراه هست.
\end{itemize}

\myopage{۵}{pic05.jpg — 20261012\_185828254.jpg}
\subsection*{تظاهرات بالینی میوکاردیت \en{chronic Active}}
\begin{itemize}
\item معمولاً بیماران مسن.
\item شروع تدریجی دارند.
\item علائم مشابه با نارسایی قلبی: خستگی، دیسپنه.
\item بیوپسی: \en{Generalized chronic Myopathic Changes, Fibrosis}.
\item در برخی بیماران به سمت نارسایی دیاستولیک می‌رود که در بیوپسی آن‌ها فیبروز وسیع و مشابه کاردیومیوپاتی رستریکتیو می‌شوند.
\item \en{60--70\%} بیماران با \en{DCM} ایدیوپاتیک را شامل می‌شود.
\item تصویربرداری‌های نوین: \en{MRI} با کنتراست گادولینیوم، \en{PET-CT}، \en{PCR}.
\end{itemize}
\subsection*{تظاهرات بالینی میوکاردیت ائوزینوفیلیک}
۳ فرم مختلف دارد:
\begin{enumerate}[label=\en{\alph*)}]
\item میوکاردیت ائوزینوفیلیک آلرژیک: اغلب آلرژی به یک دارو.
  \begin{itemize}
  \item می‌تواند شروع یک‌باره و \en{LVEF} افت کند.
  \item قطع دارو + کورتون: بهبودی و بازگشت \en{EF}.
  \end{itemize}
\item میوکاردیت در زمینهٔ اختلال ائوزینوفیلیک سیستمیک: \en{Löffler Endocarditis}.
\item میوکاردیت نکروتیک برق‌آسا:
  \begin{itemize}
  \item علت پیدا نمی‌شود: ایدیوپاتیک.
  \item سرکوب ایمنی تهاجمی.
  \end{itemize}
\end{enumerate}

\myopage{۶}{pic06.jpg — 20261012\_185834576.jpg}
\subsection*{تظاهرات بالینی \en{Peripartum Cardiomyopathy}}
\begin{itemize}
\item شروع نارسایی بطن چپ در آخرین ماه بارداری تا ۵ ماه پس از زایمان بدون وجود بیماری قلبی قبلی و علت شناخته‌شده برای کاردیومیوپاتی.
\item در \en{EMB} اکثر آن‌ها میوکاردیت وجود دارد.
\item با درمان استاندارد اغلب برمی‌گردند، پس به همین دلیل نیاز به بیوپسی وجود ندارد.
\end{itemize}
\subsection*{تشخیص میوکاردیت}
\subsubsection*{تست‌های آزمایشگاهی}
\begin{itemize}
\item \en{Troponin}: شدت بیماری، پروگنوز کوتاه‌مدت. \en{$\uparrow$ Trop} با \en{$\downarrow$ EF} همراه است.
\item \en{$\mathrm{CK\!-\!MB}>29.5\ \mathrm{ng/mL}$}: ریسک بالای مرگ در میوکاردیت.
\item در کودکان، میوکاردیت برق‌آسا، سطح بالاتر \en{Cr}، لاکتات، \en{AST} با خطر مرگ بالاتری همراه هستند.
\item در کودکان با \en{DCM} حاد در زمینهٔ میوکاردیت \en{NT-pro-BNP} افزایش پیدا می‌کند و به دنبال بهبودی آن کاهش می‌یابد.
\end{itemize}

\myopage{۷}{pic07.jpg — 20261012\_185841733.jpg}
\begin{itemize}
\item در بزرگسالان سطح بالاتر \en{IL-10} با خطر بالاتر مرگ همراه است.
\item مارکرهای التهابی مثل \en{WBC}، \en{CRP} و \en{ESR} اختصاصیت پایین دارند.
\item بررسی آنتی‌ژن‌های ویروسی در گردش خون کمکی نمی‌کند.
\item در \en{EKG}: یافتهٔ پاتوگنومونیک نداریم.
  \begin{itemize}
  \item سینوس تاکی‌کاردی شایع‌ترین یافته است.
  \item تغییرات غیراختصاصی رپولاریزاسیون.
  \item \en{PR depression}.
  \item \en{STE} منتشر.
  \end{itemize}
\end{itemize}
\subsection*{تصویربرداری‌های قلبی}
\begin{itemize}
\item ارزیابی عملکرد \en{LV} در هر بیمار با شک به میوکاردیت ضروری است.
\item رژیونال دیلاتاسیون اتاقک شایع است.
\item نارسایی بطن راست ناشایع‌تر است و اکو برای تخمین پروگنوز مهم.
\item وجود پری‌کاردیال افیوژن: میوپریکاردیت.
\item \en{WMA} معمولاً وجود دارد.
\item \en{CMR}: تست اول برای بررسی التهاب میوکارد در شک به میوکاردیت.
\end{itemize}

\myopage{۸}{pic08.jpg — 20261012\_185849105.jpg}
\begin{itemize}
\item نکروز میوکارد = \en{Late Gadolinium Enhancement (LGE)} در \en{CMR}.
\item تشخیص کل آسیب: \en{$T_1$-Weighted}.
\item تشخیص ادم میوکارد: \en{$T_2$-Weighted}.
\item \en{CMR} پس از ۲ هفته از شروع علائم دیگر تقریباً ارزش ندارد.
\item بهترین زمان در فاز حاد که \en{CP + $\uparrow$ Trop} وجود دارد.
\item بیماران با کاردیومیوپاتی یا آریتمی بطنی یا نارسایی پیشروندهٔ قلبی از درمان با استروئید سود می‌برند.
\item \en{EMB}: ریسک عارضهٔ جدی به دنبال \en{EMB}، ۱ در هر ۱۰۰۰ مورد است.
\item بیشترین جایی که \en{EMB} کمک می‌کند:
  \begin{itemize}
  \item شک به \en{GCM}، میوکاردیت ائوزینوفیلیک برق‌آسا در بیمار با کاردیومیوپاتی حاد.
  \end{itemize}
\item در بیمار با \en{Acute DCM} که به درمان جواب نمی‌دهد یا دچار بلاک یا \en{VT} پایدار می‌شود، باید به \en{GCM} شک کرد.
\item درمان \en{GCM}: سرکوب ایمنی با رژیم حاوی سیکلوسپورین.
\end{itemize}

\myopage{۹}{pic09.jpg — 20261012\_185856763.jpg}
\subsection*{درمان}
\begin{itemize}
\item خط اول درمان برای همهٔ بیماران درمان حمایتی است.
\item درصدی از بیماران به حمایت همودینامیکی نیاز پیدا می‌کنند: از وازوپرسور تا اینترا آئورتیک بالون پمپ.
\item در بیمار با \en{Acute DCM} و سندرم \en{HF}، درمان باید بر اساس گایدلاین \en{AHA} برای این \en{HF} باشد.
\item درمان روتین با سرکوب‌کننده‌های ایمنی در میوکاردیت حاد خفیف تا \en{moderately severe} در بزرگسالان توصیه نمی‌شود.
\item در بیمار \en{DCM} مزمن که به درمان \en{HF} جواب نمی‌دهد، درمان سرکوب ایمنی \en{short-term} می‌تواند کمک‌کننده باشد.
\item بیمار با بلاک قلبی یا آریتمی بطنی در زمینهٔ میوکاردیت حاد باید بستری شود.
\item آریتمی عموماً ظرف چند هفته رفع می‌شود.
\item \en{ICD} را بر اساس اندیکاسیون‌های \en{DCM} می‌گذاریم.
\end{itemize}
\subsubsection*{نمودار درمان}
% All labels and directed transitions of the handwritten flowchart are below.
\begin{latin}
\begin{description}[style=nextline,leftmargin=1em]
\item[Pts with myocarditis]
$\longrightarrow$ stable hemodynamics (Diuretics, Vasodilators).
\item[stable hemodynamics]
Stable $\longrightarrow$ Remodeling Therapy (ACEI/ARB, $\beta$ blocker $\pm$ Aldo Antagonist).\par
Unstable $\longrightarrow$ Hemodynamic support (inotrop, ballon pump).
\item[Hemodynamic support]
Stable $\longrightarrow$ Remodeling Therapy.\par
Unstable $\longrightarrow$ Immune Therapy: Bx, Steroids/Azathioprine, Interferone.
\item[Remodeling Therapy]
Stable $\longrightarrow$ Follow up: Echo, CMR, Continue therapy; if $\mathrm{LVEF}<35\%\longrightarrow$ AICD.
\item[Immune Therapy]
Unstable $\longrightarrow$ Mechanical Assist, Cardiac transplantation.
\item[Mechanical Assist, Cardiac transplantation]
$\longrightarrow$ Follow up.
\end{description}
\end{latin}

\myopage{۱۰}{pic10.jpg — 20261012\_185901813.jpg}
\begin{itemize}
\item ۱۸\% بیماران با میوکاردیت، دچار میوکاردیت راجعه می‌شوند.
\item ریسک فاکتورهای مرگ در بیماران با \en{Acute DCM}: \en{$\downarrow$ LVEF, $\downarrow$ RVEF, $\uparrow$ PAP}.
\item در کودکان می‌تواند تا ۸ سال بازگشت \en{LVEF} زمان ببرد.
\end{itemize}
\begin{latin}
\subsection*{Unexplained Acute Cardiomyopathy}
\begin{itemize}
\item Requiring Intropic Support
\item Mobitz type 2 or Higher HB
\item Sustained or Symptomatic VT
\item Failure to response to GDMT within 1--2 weeks
\end{itemize}
$\Rightarrow$ IF NO $\Rightarrow$ CMR, COR 2B\par
$\Rightarrow$ IF YES $\Rightarrow$ EMB, COR 1
\end{latin}
% Each source page repeats the subject and the blank date field, reproduced above.
% Underlined source page numbers 1--10 are represented in the page headings.
% The small printed notebook icon contains no additional text.
\ifnum\MyocarditisStandalone=1
  \let\MyocarditisEnd\enddocument
\else
  \let\MyocarditisEnd\relax
\fi
\MyocarditisEnd
